\chapter{Evaluation and Results}
\label{ch:eval}

This chapter evaluates the proposed guidance method and presents the
corresponding experimental results. The analysis is organized around the
state change induced by guidance relative to a matched unguided realization.

We first use descriptive and qualitative diagnostics to establish how guidance
modifies the generated state and how the intervention develops over flow time
(T0). We then compare alternative flow-time guidance profiles from three
complementary perspectives (T1): propagation across atmospheric variables and
vertical levels, spatial deviation from the prescribed target mask, and
diversity across stochastic ensemble realizations. Finally, after selecting a
guidance profile, we examine how the generated response changes with the
prescribed target intensity, location, and spatial extent (T2).


\section{Evaluation setup}
\label{sec:evaluation-setup}

All experiments considered in T0 and T1 generate a single weather step. This
isolates the behaviour of the guidance procedure from effects introduced by
autoregressive propagation over subsequent forecast steps.


\subsection{Experimental settings}
\label{sec:evaluation-settings}

We denote an experimental setting by

\begin{equation}
    e=(\tau,\rho,\pi),
    \label{eq:evaluation-setting}
\end{equation}

where $\tau$ is the forecast initialization, $\rho$ the prescribed relative
target intensity, and $\pi$ the spatial target mask defined in
\Cref{sec:target-region}. (TODO: not so clear) The set of evaluated settings is denoted by $E$.

A flow-time guidance profile $\gamma$ specifies how the relative guidance
strength is distributed over the residual-flow steps. We denote the set of
profiles compared within an experiment by

\begin{equation}
    \Gamma
    =
    \left\{
        \gamma^{(1)},\ldots,\gamma^{(K)}
    \right\}.
    \label{eq:candidate-guidance-profiles}
\end{equation}

The scalar guidance gain is adapted independently for each ensemble member to
realize its prescribed target, while $\gamma$ determines the temporal
distribution of the intervention.

Throughout the chapter, the variable and level conventions from
\Cref{sec:archesweather-dataset} are used. In particular,
$\ell\in L_v$ denotes a level on which variable $v$ is defined.

For quantities restricted to the spatial support of the target mask, we write

\begin{equation}
    \mathcal{R}
    :=
    \left\{
        (i,j):\pi_{ij}>0
    \right\}.
    \label{eq:target-region}
\end{equation}


\subsection{Paired stochastic realizations}
\label{sec:experimental-seeding}

Comparisons between guidance configurations use matched stochastic
realizations. The random seed is determined by the weather-step and
ensemble-member indices as

\begin{equation}
    \mathrm{seed}(n,m)=1000n+m.
    \label{eq:experimental-seed}
\end{equation}

The same seed is retained across configurations that are compared directly.
This pairs their initial residual-noise realizations and (TODO: specify more directly) removes sampling
variability in the comparison while preserving stochastic differences between
ensemble members.


\subsection{Guidance effect}
\label{sec:guidance-effect}

The central object throughout the evaluation is the state change attributable
to guidance. For experimental setting $e$, ensemble member $m$, and guidance
profile $\gamma$, we define

(TODO: do not state the numbers but the sets defined by the Dataset section)
\begin{equation}
    \Delta x_{e,m}^{\mathrm{GE},\gamma}
    :=
    x_{e,m}^{\mathrm{gui},\gamma}
    -
    x_{e,m}^{\mathrm{ung}\mid\mathrm{gui},\gamma}
    \in
    \mathbb{R}^{82\times121\times240}.
    \label{eq:guidance-effect}
\end{equation}

The guided state and its unguided-under-guided-context counterpart share the
same autoregressive conditioning context, deterministic forecast, and initial
residual-noise realization. They differ only in whether guidance is applied
during the residual flow. Their difference therefore isolates the state change
induced by guidance.

To keep the notation compact, indices for the experimental setting, ensemble
member, and guidance profile are omitted whenever they are fixed by context.
For example,

\[
    \Delta x_{v\ell}^{\mathrm{GE}}
\]

denotes the guidance-effect field for variable $v$ and level
$\ell\in L_v$ under the currently considered setting, member, and profile.

(TODO: not so clear)
The three T1 criteria examine this same object from different perspectives:
its point-wise magnitude
$|\Delta x^{\mathrm{GE}}|$, the spatial distribution of that magnitude within
the target region, and the ensemble variability of the signed effect.

\subsection{Experimental configuration}
\label{sec:experimental-configuration}

The T0 and T1 experiments target 2\,m temperature using the fixed European
mask $\pi_{\mathrm{EU}}$ introduced in \Cref{fig:target-mask-example}. All
forecast initializations are taken at 12:00 UTC. We consider the four dates
(TODO: should state explicitly that the left one is left from the peak and the other right from the peak)
\[
\begin{aligned}
    &\text{2024-08-09}, \qquad
    \text{2024-08-12}, \\
    &\text{2026-07-04}, \qquad
    \text{2026-07-09},
\end{aligned}
\]

and the two relative target intensities

\begin{equation}
    \rho \in \{0.002,\,0.004\},
    \label{eq:t1-target-intensities}
\end{equation}

corresponding to increases of $0.2\%$ and $0.4\%$ relative to each member's
unguided reference value.

(TODO: is this even necessary to say? e is e irrespectively ...)
Since the target mask is fixed throughout T1, we abbreviate the experimental
setting as

\[
    e=(\tau,\rho).
\]

The four initializations and two intensities yield eight experimental settings.
For each setting, we generate $M=5$ stochastic ensemble members. A
member-specific criterion is therefore evaluated on $8\times5=40$ paired
realizations.

% not necessary in this form: All experiments use the same numerical sampling and guidance implementation.
The residual flow is discretized into 25 flow steps.

% TODO: Add secant-search tolerance and maximum number of iterations.
% TODO: Add numerical precision used during inference.
% TODO: Add compute hardware.


\section{T0: Descriptive and qualitative evaluation}
\label{sec:qualitative-evaluation}
(TODO: throughout the descriptions in the captions should be kept to the minimum)

We begin by inspecting the generated intervention directly before reducing it
to the quantitative criteria used in T1. Unless stated otherwise, the examples
in this section use

(TODO: do we need to define e\_ex? just set the value := ?)
\begin{equation}
    e_{\mathrm{ex}}
    =
    (\text{2026-07-09},\,0.002),
    \qquad
    m=0,
    \qquad
    \gamma=\mathrm{spike@1}.
    \label{eq:representative-experiment}
\end{equation}

The profile $\mathrm{spike@1}$ is compared systematically with the alternative
flow-time profiles in T1.


\subsection{T0.1: Descriptive statistics}
\label{sec:t0-descriptive-statistics}

Add the missing two +0.06 and +0.08

We first place the prescribed changes in physical context. For each
initialization and intensity, we report the minimum and maximum 2\,m
temperature over the target support $\mathcal{R}$, together with its
area-weighted mean $\mathcal{M}_{\pi_{\mathrm{EU}}}$ and corresponding
weighted spatial standard deviation. The statistics are shown for both the
guided state and its matched unguided counterpart.

(The table is clear, we should shorten the description to the minimum)
\begin{table}[H]
    \centering
    \caption{
        Descriptive statistics of 2\,m temperature for $\rho=0.002$.
        Minimum and maximum are evaluated over the target support, while mean
        and spatial standard deviation use the area-weighted mask
        $\pi_{\mathrm{EU}}$. Entries report mean $\pm$ population standard
        deviation across the evaluated realizations, in $^\circ$C.
    }
    \label{tab:t0-descriptive-02}
    \resizebox{\textwidth}{!}{
    \begin{tabular}{lcccccccc}
        \toprule
        & \multicolumn{2}{c}{Minimum}
        & \multicolumn{2}{c}{Maximum}
        & \multicolumn{2}{c}{Mean}
        & \multicolumn{2}{c}{Spatial std.} \\
        \cmidrule(lr){2-3}
        \cmidrule(lr){4-5}
        \cmidrule(lr){6-7}
        \cmidrule(lr){8-9}
        Initialization
        & Unguided & Guided
        & Unguided & Guided
        & Unguided & Guided
        & Unguided & Guided \\
        \midrule
        2024-08-09
        & $13.03\pm0.12$ & $13.19\pm0.16$
        & $37.21\pm0.31$ & $37.67\pm0.34$
        & $24.83\pm0.046$ & $25.43\pm0.046$
        & $5.606\pm0.035$ & $5.672\pm0.033$ \\

        2024-08-12
        & $13.39\pm0.32$ & $13.57\pm0.33$
        & $38.09\pm0.40$ & $38.54\pm0.41$
        & $25.14\pm0.068$ & $25.73\pm0.069$
        & $6.024\pm0.036$ & $6.094\pm0.044$ \\

        2026-07-04
        & $12.04\pm0.12$ & $12.30\pm0.16$
        & $37.24\pm0.47$ & $38.00\pm0.49$
        & $22.66\pm0.039$ & $23.25\pm0.039$
        & $5.437\pm0.042$ & $5.436\pm0.058$ \\

        2026-07-09
        & $12.38\pm0.11$ & $12.57\pm0.14$
        & $35.51\pm0.30$ & $36.75\pm0.57$
        & $23.86\pm0.053$ & $24.46\pm0.053$
        & $5.275\pm0.031$ & $5.354\pm0.048$ \\
        \bottomrule
    \end{tabular}
    }
\end{table}

\begin{table}[H]
    \centering
    \caption{
        Descriptive statistics of 2\,m temperature for $\rho=0.004$, using
        the same convention as \Cref{tab:t0-descriptive-02}.
    }
    \label{tab:t0-descriptive-04}
    \resizebox{\textwidth}{!}{
    \begin{tabular}{lcccccccc}
        \toprule
        & \multicolumn{2}{c}{Minimum}
        & \multicolumn{2}{c}{Maximum}
        & \multicolumn{2}{c}{Mean}
        & \multicolumn{2}{c}{Spatial std.} \\
        \cmidrule(lr){2-3}
        \cmidrule(lr){4-5}
        \cmidrule(lr){6-7}
        \cmidrule(lr){8-9}
        Initialization
        & Unguided & Guided
        & Unguided & Guided
        & Unguided & Guided
        & Unguided & Guided \\
        \midrule
        2024-08-09
        & $13.03\pm0.12$ & $13.33\pm0.25$
        & $37.21\pm0.31$ & $38.13\pm0.43$
        & $24.83\pm0.046$ & $26.02\pm0.047$
        & $5.606\pm0.035$ & $5.764\pm0.038$ \\

        2024-08-12
        & $13.39\pm0.32$ & $13.75\pm0.39$
        & $38.09\pm0.40$ & $39.00\pm0.45$
        & $25.14\pm0.068$ & $26.33\pm0.069$
        & $6.024\pm0.036$ & $6.200\pm0.094$ \\

        2026-07-04
        & $12.04\pm0.12$ & $12.54\pm0.24$
        & $37.24\pm0.47$ & $38.78\pm0.54$
        & $22.66\pm0.039$ & $23.84\pm0.039$
        & $5.437\pm0.042$ & $5.453\pm0.083$ \\

        2026-07-09
        & $12.38\pm0.11$ & $12.71\pm0.21$
        & $35.51\pm0.30$ & $38.05\pm0.99$
        & $23.86\pm0.053$ & $25.05\pm0.053$
        & $5.275\pm0.031$ & $5.445\pm0.092$ \\
        \bottomrule
    \end{tabular}
    }
\end{table}

The prescribed relative changes translate into regional-mean temperature
increases of approximately $0.6^\circ$C and $1.2^\circ$C for
$\rho=0.002$ and $\rho=0.004$, respectively. Changes in the extrema and
spatial standard deviation are less uniform, indicating that realizing the
regional-mean target does not amount to adding a spatially constant
temperature offset.


\subsection{T0.2: Guided and matched unguided states}
\label{sec:t0-state-comparison}

We next inspect the three objects on which the subsequent evaluation is built.
For the representative experiment,
\Cref{fig:t0-state-comparison} shows the 2\,m temperature field of the matched
unguided state, the guided state, and their difference
$\Delta x^{\mathrm{GE}}$.

\begin{figure}[H]
    \centering

    \includegraphics[width=0.32\textwidth]{figures/T02a.png}
    \hfill
    \includegraphics[width=0.32\textwidth]{figures/T02b.png}
    \hfill
    \includegraphics[width=0.32\textwidth]{figures/T02c.png}

    \caption{
        Matched unguided state
        $x^{\mathrm{ung}\mid\mathrm{gui}}$ (left), guided state
        $x^{\mathrm{gui}}$ (center), and guidance effect
        $\Delta x^{\mathrm{GE}}$ (right) for 2\,m temperature in the
        representative experiment. The two state panels share a common scale;
        the guidance effect is shown on a separate diverging scale centered at
        zero.
    }
    \label{fig:t0-state-comparison}
\end{figure}

The guidance-effect panel separates the intervention from the weather pattern
already present in the matched unguided realization and therefore provides the
most direct view of what guidance has changed.


\subsection{T0.3: Guidance effect across atmospheric variables}
\label{sec:t0-cross-variable-effect}

Although guidance targets 2\,m temperature, its effect need not remain confined
to the target channel. For the same experimental setting and ensemble member,
\Cref{fig:t0-cross-variable-effect} shows
$\Delta x^{\mathrm{GE}}$ for the atmospheric variables used in the subsequent
flow-time analysis. Surface variables are shown at $\ell=0$ and upper-air
variables at $1000$\,hPa.

\begin{figure}[H]
    \centering

    \includegraphics[width=0.48\textwidth]{figures/T03a.png}
    \hfill
    \includegraphics[width=0.48\textwidth]{figures/T03b.png}

    \includegraphics[width=0.48\textwidth]{figures/T03c.png}
    \hfill
    \includegraphics[width=0.48\textwidth]{figures/T03d.png}

    \includegraphics[width=0.48\textwidth]{figures/T03e.png}
    \hfill
    \includegraphics[width=0.48\textwidth]{figures/T03f.png}

    \includegraphics[width=0.48\textwidth]{figures/T03g.png}

    \caption{
        Guidance effect $\Delta x^{\mathrm{GE}}$ for selected atmospheric
        variables in the representative experiment. Surface variables are
        shown at $\ell=0$ and upper-air variables at $1000$\,hPa. The maps
        provide a qualitative view of how an intervention targeted at 2\,m
        temperature propagates to other components of the atmospheric state.
    }
    \label{fig:t0-cross-variable-effect}
\end{figure}

The non-zero response outside the targeted temperature channel motivates the
multivariate propagation criterion introduced in T1.1.


\subsection{T0.4: Absolute atmospheric intensity}
\label{sec:t0-absolute-intensity}

The spatial maps show the structure of the intervention but provide limited
information about its magnitude relative to the underlying atmospheric state.
We therefore compare the area-weighted masked values

\begin{equation}
    \mathcal{M}_{\pi}
    \left(
        x_{v\ell}
    \right),
    \qquad
    \ell\in L_v,
    \label{eq:t0-absolute-intensity}
\end{equation}

across variables and levels.

\begin{figure}[H]
    \centering
    \includegraphics[width=\textwidth]{figures/T04.png}
    \caption{
        Area-weighted masked atmospheric values across variables and vertical
        levels. Solid curves show the guided states for the compared guidance
        profiles; the matched unguided state and verifying ground truth are
        shown as references.
    }
    \label{fig:t0-absolute-intensity}
\end{figure}

This view anchors the guidance-induced changes to the physical magnitude and
vertical structure of the corresponding atmospheric fields.


\subsection{T0.5: Deviation from ground truth}
\label{sec:t0-relative-ground-truth}

To complement the absolute values, we also measure each variable-level field
relative to the verifying ERA5 state,

\begin{equation}
    \mathcal{M}_{\pi}
    \left(
        x_{v\ell}
    \right)
    -
    \mathcal{M}_{\pi}
    \left(
        x_{v\ell}^{\mathrm{gt}}
    \right).
    \label{eq:t0-ground-truth-deviation}
\end{equation}

\begin{figure}[H]
    \centering
    \includegraphics[width=\textwidth]{figures/T05.png}
    \caption{
        Area-weighted masked atmospheric values relative to the verifying
        ground truth across variables and vertical levels. Zero corresponds to
        the verifying state; the matched unguided realization is shown as a
        reference.
    }
    \label{fig:t0-relative-ground-truth}
\end{figure}

Together, T0.4 and T0.5 show both the absolute atmospheric state and how far
its individual components lie from the corresponding verifying values.


\subsection{T0.6: Evolution through flow time}
\label{sec:t0-flow-convergence}

We next inspect how the intervention develops during the residual flow. The
intermediate state $x_t$ is defined in \Cref{TODO_BACKGROUND_REFERENCE}; here
we only introduce the quantities used for visualization.

For the targeted variable and level, we track the difference from the
prescribed target,

\begin{equation}
    \xi_t
    :=
    \mathcal{M}_{\pi}
    \left(
        x_{t,v^\star\ell^\star}^{\mathrm{gui}}
    \right)
    -
    y^\star.
    \label{eq:intermediate-target-gap}
\end{equation}

For the remaining variable-level fields, we instead show the masked guidance
effect,

\begin{equation}
    \mathcal{M}_{\pi}
    \left(
        x_{t,v\ell}^{\mathrm{gui}}
        -
        x_{t,v\ell}^{\mathrm{ung}\mid\mathrm{gui}}
    \right).
    \label{eq:intermediate-guidance-effect}
\end{equation}

(TODO: we only need one of the two, not both)
\begin{figure}[H]
    \centering

    \includegraphics[width=0.48\textwidth]{figures/T06a.png}
    \hfill
    \includegraphics[width=0.48\textwidth]{figures/T06b.png}

    \includegraphics[width=0.48\textwidth]{figures/T06c.png}
    \hfill
    \includegraphics[width=0.48\textwidth]{figures/T06d.png}

    \caption{
        Evolution through flow time for the configurations indicated in each
        panel. For the targeted 2\,m temperature field, the plotted quantity
        is the target gap $\xi_t$; the remaining variables show the masked
        guidance effect relative to the matched unguided trajectory. Solid
        curves show the selected ensemble member and the shaded region spans
        the minimum and maximum across the five members.
    }
    \label{fig:t0-flow-convergence}
\end{figure}

The target trajectory shows how the prescribed condition is approached during
sampling, while the remaining fields reveal when the coupled atmospheric
response emerges.

(TODO: say the thing about the shading)

\subsection{T0.7: Flow and guidance landings}
\label{sec:t0-flow-decomposition}

(TODO: we should just remove q and use the mask operator itelf, it's much clearer)
Finally, we decompose the trajectories from T0.6 into the two operations
performed at each guided flow step. For the scalar quantity displayed in a
given panel, we denote the current value by $q_t$, the value obtained after
the unguided flow move by $q_t^{\mathrm{flow}}$, and the value after the
guidance correction by $q_t^{\mathrm{gui}}$. Each step can therefore be read
as

\begin{equation}
    q_t
    \longrightarrow
    q_t^{\mathrm{flow}}
    \longrightarrow
    q_t^{\mathrm{gui}},
    \label{eq:t0-flow-guidance-landings}
\end{equation}

where the first displacement is generated by the residual flow and the second
by guidance.

(TODO: We should only stack two images, and use spikeat1 against 17 instead of different dates)
\begin{figure}[H]
    \centering

    \includegraphics[width=0.48\textwidth]{figures/T07a.png}
    \hfill
    \includegraphics[width=0.48\textwidth]{figures/T07b.png}

    \includegraphics[width=0.48\textwidth]{figures/T07c.png}
    \hfill
    \includegraphics[width=0.48\textwidth]{figures/T07d.png}

    \caption{
        Step-wise decomposition of the flow-time trajectories. Each panel
        separates the unguided flow move from the subsequent guidance move,
        making the temporal support of the applied guidance profile explicit.
    }
    \label{fig:t0-flow-decomposition}
\end{figure}

T0 therefore provides complementary views of the same intervention: its
physical magnitude, spatial structure, cross-variable response, and evolution
through the residual flow. T1 turns these observations into quantitative
criteria for comparing alternative guidance profiles.


\subsection{Target-realization check}
\label{sec:target-realization-check}

Before comparing how different profiles realize a target, we verify that the
target itself is reached. Recall from \Cref{sec:target-trajectory} that the
absolute target is constructed from the independent unguided reference
$x^{\mathrm{ung}}$, whereas
$x^{\mathrm{ung}\mid\mathrm{gui}}$ serves as the matched counterfactual used
to isolate the guidance effect.


(TODO: compress and set 1\% instead of 10e-2)
For a guided realization, we define its relative target miss as

\begin{equation}
    r_{e,m}^{\gamma}
    :=
    \frac{
        \left|
            \mathcal{M}_{\pi_e}
            \left(
                x_{e,m,v^\star\ell^\star}^{\mathrm{gui},\gamma}
            \right)
            -
            y_{e,m}^{\star}
        \right|
    }{
        \left|
            \mathcal{M}_{\pi_e}
            \left(
                x_{e,m,v^\star\ell^\star}^{\mathrm{ung}}
            \right)
            - y_{e,m}^{\star}
        \right|
    }.
    \label{eq:target-relative-miss}
\end{equation}

We consider the target successfully realized when

\begin{equation}
    r_{e,m}^{\gamma}
    \leq
    \varepsilon_{\mathrm{tar}},
    \qquad
    \varepsilon_{\mathrm{tar}}=10^{-2}.
    \label{eq:target-realization-tolerance}
\end{equation}

This check is treated as a prerequisite rather than as a criterion for ranking
guidance profiles. T1 compares the atmospheric responses of profiles that
successfully realize the prescribed target.
(TODO: all our experiments converge succesfully)

\section{T1: Guidance-profile evaluation}
\label{sec:guidance-profile-evaluation}

We now compare how the atmospheric response depends on the flow-time
distribution of guidance. The comparison is performed in two stages. We first
localize the intervention at individual flow steps and then test whether
spreading guidance over neighbouring steps improves the response.

For the localization experiment, we consider

\begin{equation}
    \Gamma_{\mathrm{spike}}
    =
    \left\{
        \mathrm{spike@0},
        \mathrm{spike@1},
        \mathrm{spike@2},
        \mathrm{spike@5},
        \mathrm{spike@11},
        \mathrm{spike@17},
        \mathrm{spike@23}
    \right\},
    \label{eq:spike-profile-set}
\end{equation}

where $\mathrm{spike@}k$ assigns unit relative guidance strength to flow step
$k$ and zero to all remaining steps.

The profiles are evaluated from three complementary perspectives. T1.1
measures the magnitude of the guidance effect across atmospheric variables and
levels, T1.2 measures how its spatial structure differs from the prescribed
mask, and T1.3 measures how strongly the induced response varies across
ensemble members.


\subsection{Common score construction}
\label{sec:t1-common-score}

(TODO: actually it's probably best to use all indices)

Each criterion produces a non-negative value $c_{v\ell}^{\gamma}$ for every
variable $v$, level $\ell\in L_v$, and candidate profile $\gamma$. For T1.1
and T1.2 this value is computed separately for each experimental setting and
ensemble member; these indices are omitted here for readability.

We first aggregate over the levels of each variable and normalize relative to
the strongest candidate,

\begin{equation}
    N_v^\gamma
    :=
    \frac{
        \displaystyle
        \sum_{\ell\in L_v}
        c_{v\ell}^{\gamma}
    }{
        \displaystyle
        \max_{\gamma'\in\Gamma}
        \sum_{\ell\in L_v}
        c_{v\ell}^{\gamma'}
    }.
    \label{eq:t1-variable-score}
\end{equation}

This gives equal importance to variables whose raw values may differ strongly
in magnitude or physical units. The normalized variable scores are combined
into a single multivariate score,

\begin{equation}
    R^\gamma
    :=
    \frac{
        \displaystyle
        \sum_{v} N_v^\gamma
    }{
        \displaystyle
        \max_{\gamma'\in\Gamma}
        \sum_v N_v^{\gamma'}
    }.
    \label{eq:t1-single-score}
\end{equation}

Hence $R^\gamma\in[0,1]$, with the best candidate for a given realization
attaining one. Since the normalization is relative to the candidate set
$\Gamma$, all scores are recomputed whenever $\Gamma$ changes.

All reported spreads below are population standard deviations. The dimensions
over which values are aggregated are stated in the corresponding table or
figure caption.


\subsection{T1.1: Multivariate propagation}
\label{sec:multivariate-propagation}

The first criterion asks how strongly the intervention propagates through the
atmospheric state. For each variable and level, we evaluate the area-weighted
mean magnitude of the guidance effect within the target region,

\begin{equation}
    c_{v\ell}^{\gamma}
    :=
    \mathcal{M}_{\pi}
    \left(
        \left|
            \Delta x_{v\ell}^{\mathrm{GE},\gamma}
        \right|
    \right).
    \label{eq:t1-propagation}
\end{equation}

The absolute value prevents positive and negative spatial responses from
cancelling. Larger values therefore indicate a stronger intervention in the
corresponding atmospheric field.

Before reducing the response spatially,
\Cref{fig:t1-propagation-example} shows the signed guidance effect directly
for selected variables throughout the residual flow.

(TODO: add also specific humidity!)
\begin{figure}[H]
    \centering
    \includegraphics[width=\textwidth]{figures/T1.1a.png}
    \caption{
        Evolution of the signed guidance effect
        $\Delta x^{\mathrm{GE}}$ through flow time for selected atmospheric
        variables in a representative experiment. The maps show the
        unaggregated object underlying the multivariate-propagation criterion.
    }
    \label{fig:t1-propagation-example}
\end{figure}

The level-resolved values from \Cref{eq:t1-propagation} are shown in
\Cref{fig:t1-propagation}. For each variable, the horizontal axis traverses
its available levels. The solid curve corresponds to one ensemble member and
the band spans the minimum and maximum across all five members.

\begin{figure}[H]
    \centering
    \includegraphics[width=\textwidth]{figures/T1.1.png}
    \caption{
        Level-resolved multivariate propagation for the candidate spike
        profiles. Each value is the area-weighted masked mean
        $\mathcal{M}_{\pi}(|\Delta x^{\mathrm{GE}}_{v\ell}|)$.
        Solid curves show one ensemble member and bands span the minimum and
        maximum across the five members.
    }
    \label{fig:t1-propagation}
\end{figure}

(TODO: wrong to interprete this here, also spikeat2 is better than 0, so should specify better)
The response exhibits a pronounced maximum at $\mathrm{spike@1}$.
$\mathrm{spike@0}$ produces a weaker response, showing that the dependence is
not simply monotonic with increasingly early guidance. From step~1 onward,
however, propagation decreases strongly as the intervention is shifted toward
later flow steps.


\subsection{T1.2: Spatial deviation from the target mask}
\label{sec:spatial-mask-deviation}

The propagation criterion measures the magnitude of the response but does not
describe its spatial structure. We therefore compare the normalized spatial
distribution of
$|\Delta x_{v\ell}^{\mathrm{GE},\gamma}|$ with the prescribed mask.

For each variable and level, we define

\begin{equation}
    c_{v\ell}^{\gamma}
    :=
    D_{v\ell}^{\gamma}
    =
    \frac{1}{2}
    \sum_{(i,j)\in\mathcal{R}}
    \left|
        \frac{
            \left|
                \Delta x_{v\ell ij}^{\mathrm{GE},\gamma}
            \right|
        }{
            \displaystyle
            \sum_{(i',j')\in\mathcal{R}}
            \left|
                \Delta x_{v\ell i'j'}^{\mathrm{GE},\gamma}
            \right|
        }
        -
        \pi_{ij}
    \right|.
    \label{eq:t1-spatial-deviation}
\end{equation}

Because both spatial distributions sum to one over $\mathcal{R}$,

\[
    D_{v\ell}^{\gamma}\in[0,1].
\]

A value of zero means that the magnitude of the guidance effect follows the
prescribed mask exactly within the target region. Larger values indicate an
increasing departure from this imposed spatial pattern.

Figure~\ref{fig:t1-spatial-deviation-example} illustrates the two spatial
objects being compared before their reduction to $D_{v\ell}^{\gamma}$.

(TODO: should again inspect z1000 or specific humidity and show the artifacts and explain on the figure)
\begin{figure}[H]
    \centering
    \includegraphics[width=\textwidth]{figures/T1.2a.png}
    \caption{
        Example normalized spatial distributions used for T1.2. The
        normalized magnitude of the guidance effect is compared with the
        prescribed spatial mask for the candidate profiles.
    }
    \label{fig:t1-spatial-deviation-example}
\end{figure}

\begin{figure}[H]
    \centering
    \includegraphics[width=\textwidth]{figures/T1.2.png}
    \caption{
        Spatial deviation $D_{v\ell}^{\gamma}$ across atmospheric variables
        and vertical levels for the candidate spike profiles. Solid curves
        show one ensemble member and bands span the minimum and maximum across
        the five members.
    }
    \label{fig:t1-spatial-deviation}
\end{figure}

We interpret larger spatial deviation as evidence that the atmospheric
response is not simply reproducing the prescribed mask geometry. It is not,
however, treated as a standalone measure of realism and is considered together
with the qualitative diagnostics and the other two T1 criteria.

(TODO: again)
Unlike propagation, the spatial-deviation criterion does not select
$\mathrm{spike@1}$ as its strongest candidate. The highest aggregate score is
obtained by $\mathrm{spike@0}$, while $\mathrm{spike@1}$ and
$\mathrm{spike@2}$ remain close. The temporal dependence is also less
monotonic than for multivariate propagation.


\subsection{T1.3: Ensemble diversity}
\label{sec:guidance-footprint-variability}

The third criterion examines whether the response induced by guidance remains
dependent on the stochastic ensemble realization. For each variable and level,
we compute

\begin{equation}
    c_{v\ell}^{\gamma}
    :=
    \mathcal{M}_{\pi}
    \left(
        \operatorname{std}_{m}
        \left[
            \Delta x_{m,v\ell}^{\mathrm{GE},\gamma}
        \right]
    \right),
    \label{eq:t1-ensemble-diversity}
\end{equation}

where the standard deviation is evaluated point-wise over the $M$ ensemble
members. The signed guidance effect is used before taking the ensemble
standard deviation.

Larger values indicate stronger member-to-member differences in the induced
response and therefore greater retention of ensemble diversity under guidance.

For visualization,
\Cref{fig:t1-ensemble-diversity-example} shows the spatial distribution of the
pixel-wise ensemble standard deviation, normalized only for display.

(TODO: again)
\begin{figure}[H]
    \centering
    \includegraphics[width=\textwidth]{figures/T1.3a.png}
    \caption{
        Example spatial distribution of ensemble diversity for the candidate
        profiles. Each panel shows the pixel-wise standard deviation of the
        signed guidance effect across ensemble members, normalized for
        visualization.
    }
    \label{fig:t1-ensemble-diversity-example}
\end{figure}

\begin{figure}[H]
    \centering
    \includegraphics[width=\textwidth]{figures/T1.3.png}
    \caption{
        Ensemble diversity across atmospheric variables and vertical levels,
        measured according to \Cref{eq:t1-ensemble-diversity}. No ensemble
        band is shown because the member dimension is already consumed by the
        standard deviation.
    }
    \label{fig:t1-ensemble-diversity}
\end{figure}

The ensemble-diversity score peaks at $\mathrm{spike@1}$.
$\mathrm{spike@0}$ again produces a weaker response, while diversity decreases
strongly as guidance is shifted from step~1 toward later flow steps.


\subsection{Temporal-localization summary}
\label{sec:t1-localization-summary}

The three criteria are summarized in
\Cref{tab:t1-localization-summary}. For propagation and spatial deviation,
the values are pooled over the four initializations, two target intensities,
and five ensemble members. Ensemble diversity is instead pooled over the four
initializations and two intensities because its member dimension is consumed
by the point-wise ensemble standard deviation.

\begin{table}[H]
    \centering
    \caption{
        Temporal-localization scores across both target intensities.
        Propagation and spatial deviation report mean $\pm$ population
        standard deviation over $4\times2\times5=40$ setting-member
        realizations. Ensemble diversity is summarized over the
        $4\times2=8$ initialization-intensity settings. Larger values are
        preferred for all three criteria.
    }
    \label{tab:t1-localization-summary}
    \begin{tabular}{lccc}
        \toprule
        Profile
        & Propagation
        & Spatial deviation
        & Ensemble diversity \\
        \midrule
        $\mathrm{spike@0}$
        & $0.689\pm0.050$
        & $\mathbf{0.998\pm0.005}$
        & $0.682\pm0.027$ \\

        $\mathrm{spike@1}$
        & $\mathbf{1.000\pm0.000}$
        & $0.977\pm0.017$
        & $\mathbf{1.000\pm0.000}$ \\

        $\mathrm{spike@2}$
        & $0.770\pm0.057$
        & $0.984\pm0.012$
        & $0.772\pm0.039$ \\

        $\mathrm{spike@5}$
        & $0.463\pm0.040$
        & $0.915\pm0.018$
        & $0.402\pm0.017$ \\

        $\mathrm{spike@11}$
        & $0.134\pm0.013$
        & $0.840\pm0.030$
        & $0.115\pm0.011$ \\

        $\mathrm{spike@17}$
        & $0.046\pm0.004$
        & $0.857\pm0.042$
        & $0.041\pm0.005$ \\

        $\mathrm{spike@23}$
        & $0.028\pm0.002$
        & $0.858\pm0.042$
        & $0.011\pm0.001$ \\
        \bottomrule
    \end{tabular}
\end{table}

% The criteria therefore provide complementary rather than identical rankings.
% $\mathrm{spike@1}$ maximizes both multivariate propagation and ensemble
% diversity, whereas $\mathrm{spike@0}$ yields the largest spatial deviation.
% The differences in spatial deviation among the early profiles are relatively
% small, while their propagation and diversity scores differ substantially.
% We therefore retain $\mathrm{spike@1}$ as the localized reference profile and
% next test whether distributing its guidance signal over neighbouring flow
% steps improves the response.

Taken together, the three criteria favour guidance concentrated near the
beginning of the residual flow. Among the evaluated profiles,
$\mathrm{spike@1}$ produces the strongest multivariate propagation and
ensemble diversity while remaining close to the strongest candidates in
spatial deviation. We therefore select

\begin{equation}
    \gamma^\star = \mathrm{spike@1}
    \label{eq:selected-guidance-profile}
\end{equation}

for the subsequent experiments.

(TODO: we can say here that we also experimented with spreading the signal but the results remained fedel to these findings)


\section{T2: Scenario-specification analysis}
\label{sec:scenario-specification-analysis}

Having fixed the guidance profile to
$\gamma^\star=\mathrm{spike@1}$, we next examine how the generated response
changes with the prescribed scenario. We separately vary the target intensity,
location, and spatial extent while keeping the remaining configuration fixed.


\subsection{T2.1: Target intensity}
\label{sec:target-intensity-analysis}

We first vary the prescribed relative intensity $\rho$ while keeping the
forecast initialization and spatial target mask fixed.

% TODO: State the evaluated intensity levels.
% \[
%     \rho \in \{\ldots\}.
% \]

For each intensity, we compare the resulting guided state with its matched
unguided counterpart and inspect the corresponding guidance effect
$\Delta x^{\mathrm{GE}}$. This shows how the magnitude and spatial structure
of the required intervention change as the prescribed target is moved further
from the unguided reference.

\begin{figure}[H]
    \centering
    % TODO: Insert target-intensity figure(s).
    \caption{
        Guided states and corresponding guidance effects for increasing target
        intensity using $\gamma^\star=\mathrm{spike@1}$. All other components
        of the experimental setting are held fixed.
    }
    \label{fig:t2-target-intensity}
\end{figure}

% TODO: Add concise result interpretation.


\subsection{T2.2: Target location}
\label{sec:target-location-analysis}

We next shift the spatial target mask while retaining its size, target
variable and level, prescribed intensity, and guidance profile.

% TODO: State the evaluated target locations.

\begin{figure}[H]
    \centering
    % TODO: Insert target-location figure(s).
    \caption{
        Guided states and guidance effects for target regions placed at
        different locations. The mask extent and prescribed intensity are kept
        fixed.
    }
    \label{fig:t2-target-location}
\end{figure}

This experiment tests whether the character of the intervention depends on
the atmospheric state present at the location where the target is imposed.

% TODO: Add concise result interpretation.


\subsection{T2.3: Target size}
\label{sec:target-size-analysis}

Finally, we vary the spatial extent of the target mask while retaining its
center, target variable and level, prescribed intensity, and guidance profile.

% TODO: State the evaluated mask sizes.

\begin{figure}[H]
    \centering
    % TODO: Insert target-size figure(s).
    \caption{
        Guided states and guidance effects for target masks of increasing
        spatial extent. The mask center and prescribed intensity are held
        fixed.
    }
    \label{fig:t2-target-size}
\end{figure}

This comparison shows how the intervention changes when the same relative
target is prescribed over increasingly large spatial regions.

% TODO: Add concise result interpretation.